\documentclass[aspectratio=169]{beamer}

\usetheme{Madrid}
\usecolortheme{default}
\setbeamertemplate{navigation symbols}{}
\setbeamertemplate{footline}[frame number]

\usepackage{amsmath}
\usepackage{amssymb}
\usepackage{booktabs}
\usepackage{graphicx}
\usepackage{hyperref}

\title[Scattering likelihood]{Scattering likelihood for CIB noise}
\subtitle{Bird's-eye view of the current pipeline}
\author{}
\date{}

\newcommand{\ph}{\phi}
\newcommand{\mucib}{\mu_{\rm CIB}}
\newcommand{\Ccib}{C_{\rm CIB}}
\newcommand{\Ulike}{U_{\rm like}}
\newcommand{\Uprior}{U_{\rm prior}}

\begin{document}

\begin{frame}
  \titlepage
\end{frame}

\begin{frame}{Goal}
  We observe
  \[
    d = s + n,
  \]
  where \(s\) is the signal and \(n\) is CIB-like non-Gaussian noise.

  \vspace{0.4cm}
  The goal is to write and sample a posterior likelihood using a learned
  scattering model for the noise:
  \[
    p(s\mid d)
    \propto
    \exp[-U(s)],
    \qquad
    U(s)
    =
    \Ulike(d-s)
    +
    \Uprior(s).
  \]

  \vspace{0.3cm}
  The scattering likelihood is
  \[
    \Ulike(n)
    =
    \frac12
    \left[\ph(n)-\mucib\right]^{\top}
    \Ccib^{-1}
    \left[\ph(n)-\mucib\right].
  \]
\end{frame}

\begin{frame}{What Is Learned}
  The learned object is not an image-space Gaussian. It is a Gaussian model
  in reduced scattering-statistic space:
  \[
    \ph(n) \sim \mathcal{N}(\mucib,\Ccib).
  \]

  \vspace{0.35cm}
  Therefore, for any candidate signal \(s\), the residual
  \[
    n(s)=d-s
  \]
  is scored by whether its scattering statistics look like CIB statistics.

  \vspace{0.35cm}
  The scalar diagnostic is
  \[
    \rho(n)
    =
    \sqrt{
    \left[\ph(n)-\mucib\right]^{\top}
    \Ccib^{-1}
    \left[\ph(n)-\mucib\right]
    },
    \qquad
    \Ulike=\frac12\rho^2.
  \]
\end{frame}

\begin{frame}{Covariance Estimation Script}
  \texttt{estimate\_scattering\_covariance.py} estimates the ingredients
  \(\mucib\), \(\Ccib\), and the metadata needed to define \(\ph\).

  \vspace{0.3cm}
  Inputs:
  \[
    n_a \in \mathbb{R}^{256\times256},
    \qquad a=1,\ldots,N.
  \]

  Operations:
  \[
    y_a = \ph(n_a),
    \qquad
    \widehat{\mu}
    =
    \frac1N\sum_{a=1}^{N} y_a,
  \]
  \[
    \widehat{C}
    =
    \frac{1}{N-1}
    \sum_{a=1}^{N}
    (y_a-\widehat{\mu})(y_a-\widehat{\mu})^{\top}.
  \]

  \vspace{0.25cm}
  If too few input maps are available, the script can synthesize additional
  maps.
\end{frame}

\begin{frame}{Current CIB Product}
  Current CIB product:
  \[
    N = 3400
  \]
  patches extracted from the original $100$ maps of size $1500 \times 1500$ 

  \vspace{0.35cm}
  Reduced statistic dimension:
  \[
    d_{\rm active}=188.
  \]

  \vspace{0.35cm}
  The JSON fixes the statistic definition: preprocessing, scattering
  parameters, normalization, and active-mask convention.
\end{frame}

\begin{frame}{Likelihood Validation}
  We tested three objects under the learned CIB likelihood:

  \vspace{0.2cm}
  \begin{itemize}
    \item Gaussian draws directly in statistic space:
    \(y\sim \mathcal{N}(\mucib,\Ccib)\).
    \item Maps synthesized to match such statistic draws.
    \item Real CIB patches from the original bank.
  \end{itemize}

  \vspace{0.3cm}
  If the model were perfectly realized by the synthesized maps, the synthesized
  maps would have \(\rho\) comparable to direct statistic draws and real CIB maps.

  \vspace{0.3cm}
  For \(d_{\rm active}=188\), the typical Gaussian-null scale is
  \[
    \rho \sim \sqrt{188} \simeq 13.7.
  \]
\end{frame}

\begin{frame}{Likelihood Validation Results}
  \begin{center}
  \begin{tabular}{lrrrrr}
    \toprule
    Object & \(n\) & mean \(\rho\) & median \(\rho\) & range \(\rho\) & mean \(\Ulike\) \\
    \midrule
    Statistic draws & 7 & 12.6862 & 12.5642 & 11.7173--13.7384 & 80.6452 \\
    Synthesized maps & 7 & 34.7629 & 33.9045 & 31.4099--43.2775 & 611.2964 \\
    Real CIB maps & 7 & 13.7758 & 13.4523 & 12.1901--17.7101 & 96.3366 \\
    \bottomrule
  \end{tabular}
  \end{center}

  \vspace{0.35cm}
  Interpretation:
  direct statistic draws and real CIB maps are on the expected scale
  \(\rho\sim\sqrt{d}\). Synthesized maps are much worse under the same
  likelihood.

  \vspace{0.25cm}
  This indicates that the map synthesis step does not faithfully realize the
  sampled scattering statistics, even when the target statistic itself is
  likely under the learned Gaussian.
\end{frame}

\begin{frame}{Sampled-Statistic Syntheses}
  \begin{center}
    \includegraphics[width=0.95\linewidth,height=0.78\textheight,keepaspectratio]{figures/validation_autoscale.png}
  \end{center}

  \vspace{-0.2cm}
  \centering
  Same validation image with individual plotting scale.
\end{frame}

\begin{frame}{Sampled-Statistic Syntheses: Fixed Scale}
  \begin{center}
    \includegraphics[width=0.95\linewidth,height=0.78\textheight,keepaspectratio]{figures/validation_fixed_scale.png}
  \end{center}

  \vspace{-0.2cm}
  \centering
  Same objects on a common color scale.
\end{frame}

\begin{frame}{MGVI Formulation}
  For posterior sampling, we used NIFTy/MGVI with
  \[
    U(s)
    =
    \frac12
    \left[\ph(d-s)-\mucib\right]^{\top}
    \Ccib^{-1}
    \left[\ph(d-s)-\mucib\right]
    +
    \Uprior(s).
  \]

  \vspace{0.3cm}
  The prior was a correlated-field Gaussian prior for \(s\). The likelihood
  is the learned scattering model for the residual \(d-s\), not for \(s\).

  \vspace{0.3cm}
  A (refined) recovered signal from a vanilla ST component separation run was used as an initialization.
 
\end{frame}

\begin{frame}{MGVI Diagnostic}
  \[
    d-s_{\rm sample}
  \]
  should look CIB-like and have reasonable scattering likelihood.

  \vspace{0.3cm}
  Latest run diagnostics:
  \[
    \Ulike(s=0)=13620.6,
    \qquad
    \Ulike(s_{\rm MAP})=39.36.
  \]

  \vspace{0.2cm}
  The MAP strongly improves the residual likelihood. The remaining issue is
  morphological: posterior samples can contain negative compact structures in
  \(s\), suggesting that the signal prior is not restrictive enough where the
  CIB likelihood allows point-source-like residual tradeoffs.
\end{frame}

\begin{frame}{MGVI Results}
  \begin{center}
    \includegraphics[width=0.98\linewidth,height=0.80\textheight,keepaspectratio]{figures/mgvi_results_signal.png}
  \end{center}
\end{frame}

\section{Direct Langevin sampling}

\begin{frame}{Direct Image-Space Langevin Sampling}
  We replace the NIFTy/MGVI sampling stage by unadjusted Langevin dynamics
  (ULA) directly in image space:
  \[
    s_{k+1}
    =s_k-\epsilon\,\nabla_s U(s_k)
    +\sqrt{2\epsilon}\,\xi_k,
    \qquad
    \xi_k\sim\mathcal{N}(0,I_{256^2}).
  \]

  At every iteration:
  \begin{enumerate}
    \item form the candidate CIB residual \(n_k=d-s_k\);
    \item compute the differentiable STL vector \(\ph(n_k)\);
    \item evaluate \(\Ulike(n_k)\) and the signal regularization;
    \item use PyTorch automatic differentiation for \(\nabla_s U\);
    \item take the deterministic drift plus Gaussian-noise step.
  \end{enumerate}

  The diagnostic
  \[
    \rho_k=\sqrt{2\Ulike(d-s_k)}
  \]
  contains only the scattering-likelihood contribution.
\end{frame}

\begin{frame}{Why a Signal-Space Constraint Is Needed}
  The image has \(256^2=65\,536\) pixel degrees of freedom, whereas the
  active likelihood vector has only \(188\) statistics. Locally,
  \[
    \delta\ph\simeq J_\ph\,\delta s,
    \qquad
    J_\ph\in\mathbb{R}^{188\times65\,536}.
  \]
  Rank--nullity therefore implies
  \[
    \dim\ker J_\ph
    =65\,536-\operatorname{rank}(J_\ph)
    \ge65\,536-188=65\,348.
  \]

  Many local image perturbations are consequently invisible to the
  likelihood at first order. Without a restoring term they undergo Langevin
  diffusion; with a prior, their behaviour is largely controlled by that
  prior.
\end{frame}

\begin{frame}{Initialization: Component Separation and Refinement}
  We do \emph{not} initialize the chain from zero or from the observed map.
  The initialization is built in two deterministic stages:
  \begin{enumerate}
    \item STL component separation gives \(s_{\rm compsep}\).
    \item We refine it by minimizing
      \[
        \Ulike(d-s)
        +\frac{\lambda_{\rm ref}}{2}
         \lVert s-s_{\rm compsep}\rVert^2,
        \qquad \lambda_{\rm ref}=0.1,
      \]
      using L-BFGS (maximum 300 iterations).
  \end{enumerate}

  The refinement changed
  \[
    \rho:64.59\longrightarrow6.82,
    \qquad
    {\rm RMS}(s_{\rm refined}-s_{\rm compsep})\simeq0.137.
  \]

  \alert{Both Langevin approaches below start from}
  \[
    s_0=s_{\rm refined},
  \]
  after imposing the zero-mean convention.
\end{frame}

\begin{frame}{Approach 1: Power-Law Gaussian Prior}
  The first ULA target was
  \[
    U(s)=\Ulike(d-s)+\frac12
    \sum_{k\ne0}\frac{|\widehat{s}(k)|^2}{P(k)},
    \qquad
    P(k)\propto k^{-3.07751}.
  \]
  This prior is zero-centred. The exponent \(-3.07751\) is the
  \emph{power-spectrum} slope, twice the Fourier-amplitude slope
  \(-1.53876\). The spectrum was normalized to target RMS \(2.53272\).

  \vspace{0.2cm}
  Run parameters:
  \[
    s_0=s_{\rm refined},\quad
    \epsilon=10^{-4},\quad
    N_{\rm step}=20\,000,\quad
    N_{\rm burn}=10\,000,
  \]
  with random seed \(240829\), float64 arithmetic, and a zero-mean signal.
\end{frame}

\begin{frame}{Approach 1: Observed Behaviour}
  Although the refined initialization had \(\rho=6.82\), the power-law prior
  assigned it a very large energy and drove the residual back toward
  \(\rho\simeq40\)--47.

  \begin{itemize}
    \item The chain correctly loaded \(s_{\rm refined}\).
    \item \(\Ulike\) increased while the power-law prior energy decreased by a
      much larger amount.
    \item Total energy therefore decreased as the residual became less
      CIB-like.
    \item Changing only the initialization cannot fix a mismatch in the
      target posterior.
  \end{itemize}

  This simple prior fixes a global RMS and spectral slope, but it does not
  encode the morphology of the successfully separated signal.
\end{frame}

\begin{frame}{Approach 2: Matching the Refined Map}
  We replace the zero-centred power-law prior by a local quadratic anchor:
  \[
    U(s)=\Ulike(d-s)
    +\frac{\lambda}{2}\lVert s-s_{\rm refined}\rVert^2.
  \]
  The initial condition and anchor centre are identical:
  \[
    s_0=s_{\rm anchor}=s_{\rm refined}.
  \]

  This gives every image-space direction restoring curvature \(\lambda\),
  while preserving the refined reconstruction. In an anchor-only direction,
  the nominal pixel width is
  \[
    \sigma_{\rm anchor}=\lambda^{-1/2}.
  \]

  We explored
  \[
    \lambda=100,50,25
    \quad\Longleftrightarrow\quad
    \sigma_{\rm anchor}=0.10,0.141,0.20.
  \]
  With \(\lambda=100\), the chain stabilized near \(\rho\simeq8.5\)--9;
  smaller \(\lambda\) increases map variability and residual \(\rho\).
\end{frame}

\begin{frame}{Current Anchored-ULA Configuration}
  \begin{center}
  \begin{tabular}{ll}
    \toprule
    Quantity & Value \\
    \midrule
    Initial map & \(s_{\rm refined}\) \\
    Anchor centre & \(s_{\rm refined}\) \\
    Anchor precision & \(\lambda=50\) (\(25\) and \(100\) also tested) \\
    ULA step size & \(\epsilon=10^{-4}\) \\
    Total steps & \(20\,000\) \\
    Burn-in & \(10\,000\) \\
    Thinning & \(1\,000\) \\
    Saved samples & \(10\) \\
    Random seed & \(240829\) \\
    Arithmetic & float64, zero-mean signal \\
    \bottomrule
  \end{tabular}
  \end{center}

  Diagnostics saved at every step are total, likelihood and anchor energies,
  gradient RMS, \(\rho\), and RMS distance from the anchor. Samples, residuals,
  posterior mean, and posterior standard deviation are also saved.
\end{frame}

\begin{frame}{Current Interpretation}
  Main conclusions:

  \vspace{0.2cm}
  \begin{itemize}
    \item The learned CIB likelihood is well-defined by \(\mucib\), \(\Ccib\),
    and the JSON scattering metadata.
    \item Real CIB maps score near the expected Gaussian-null scale in
    statistic space.
    \item Synthesizing maps from sampled statistics is currently imperfect:
    those maps have too-large \(\rho\).
    \item In MGVI, the residual \(d-s\) can look CIB-like, but the recovered
    \(s\) can still absorb negative point-source-like structures.
    \item Direct ULA confirms that the signal-space constraint is decisive:
    the zero-centred power-law prior rejects the refined solution, whereas a
    refined-map anchor produces controlled local samples.
  \end{itemize}

  \vspace{0.35cm}
  The next technical pressure point is therefore the signal model/prior, not
  only the CIB likelihood.
\end{frame}

\end{document}
